\documentclass{article}
\usepackage{amsmath,amssymb}
\usepackage{xurl}
\usepackage[hidelinks]{hyperref}
% Shared commands for notes in docs/paper-gaps/.

\DeclareUnicodeCharacter{2080}{\ifmmode{}_{0}\else\textsubscript{0}\fi}
\DeclareUnicodeCharacter{2081}{\ifmmode{}_{1}\else\textsubscript{1}\fi}
\DeclareUnicodeCharacter{2082}{\ifmmode{}_{2}\else\textsubscript{2}\fi}
\DeclareUnicodeCharacter{2099}{\ifmmode{}_{n}\else\textsubscript{n}\fi}

\newcommand{\Lean}{\textsc{Lean}}
\ExplSyntaxOn
\NewDocumentCommand{\leanid}{m}{
  \ifmmode
    \textsf{\detokenize{#1}}
  \else
    \group_begin:
    \urlstyle{sf}
    \tl_set:Nn \l_tmpa_tl {#1}
    \tl_replace_all:Nnn \l_tmpa_tl {\_} {_}
    \exp_args:NV \path \l_tmpa_tl
    \group_end:
  \fi
}
\ExplSyntaxOff
\providecommand{\tr}{\operatorname{tr}}
\newcommand{\tnleanIssueBase}{https://github.com/LionSR/TNLean/issues/}
\newcommand{\tnleanPrBase}{https://github.com/LionSR/TNLean/pull/}
\newcommand{\tnleanDocBase}{https://sirui-lu.com/TNLean/docs/}
\newcommand{\ghissue}[1]{\href{\tnleanIssueBase#1}{GitHub issue \##1}}
\newcommand{\ghpr}[1]{\href{\tnleanPrBase#1}{PR \##1}}
\newcommand{\doclink}[1]{\href{\tnleanDocBase#1.html}{\path{#1}}}

% Dirac notation and the identity operator, as in blueprint/src/macros/common.tex.
% \providecommand keeps note-local definitions authoritative.
\usepackage{dsfont}
\providecommand{\ket}[1]{|#1\rangle}
\providecommand{\bra}[1]{\langle #1|}
\providecommand{\braket}[2]{\langle #1 | #2 \rangle}
\providecommand{\ketbra}[2]{|#1\rangle\!\langle #2|}
\providecommand{\Id}{\mathds{1}}
\providecommand{\one}{\mathds{1}}
\providecommand{\C}{\mathbb{C}}
\providecommand{\MN}[1]{M_{#1}(\C)}
\providecommand{\MD}{M_D(\C)}

% Verdict marker: \gapnote{<kind>}{<status>}. Kind is the policy
% classification (clarification, local-correction, scope-restriction,
% unfaithful, false-source, open-gap); status is open, wip (elimination
% underway), resolved, or historical. Machine-read by the site index;
% prints a verdict line.
\newcommand{\gapnote}[2]{\par\noindent\textbf{Verdict:} #1 (#2).\par}

\setlength{\emergencystretch}{2em}
\title{The Finite Periodic One-Species PVBS Gap}
\author{}
\date{2026-10-04}
\begin{document}
\maketitle
\gapnote{scope-restriction}{open}

\section{Source and scope}
Bachmann and Nachtergaele, arXiv:1112.4097, Section II after equation (8),
give the one-particle dispersion
$\epsilon_i(k)=1-2\lambda_i\cos(k+\theta_{i0})/(1+\lambda_i^2)$.
They then conjecture that the exact thermodynamic gap for the multi-species
model is the minimum of $(1-\lambda_i)^2/(1+\lambda_i^2)$ over species.
The result proved here is the finite periodic one-species, zero-phase case,
with positive hopping $\mu\ne1$. It is not a formal proof of that general
thermodynamic conjecture.

\section{Normalization and the finite ring}
The local parent interaction is exactly
\begin{align}
    h_\mu & =|11\rangle\langle11|+
    \frac{(|01\rangle-\mu|10\rangle)
        (\langle01|-\mu\langle10|)}{1+\mu^2}.
    \label{eq:bn12_gap_local}
\end{align}
On a ring of $N\ge2$ sites, sum this interaction over the $N$ ordered
windows $(i,i+1)$ modulo $N$. For $N=2$, this includes both orientations;
a convention with only one undirected bond is not covered by the claimed
constant. The local formula is identified with the repository's actual
orthogonal-projector parent interaction, rather than defining a replacement
Hamiltonian.

Let $\mathcal N$ count occupied sites. The exact periodic identity is
\begin{align}
    H_\mu & =a(\mu)H_1+c(\mu)\mathcal N,
          & a(\mu)                       & =\frac{2\mu}{1+\mu^2},
          & c(\mu)                       & =\frac{(\mu-1)^2}{1+\mu^2}.
    \label{eq:bn12_gap_decomposition}
\end{align}
Indeed, the hopping coefficients agree, while the difference in diagonal
terms is $[(\mu^2-\mu)n_L+(1-\mu)n_R]/(1+\mu^2)$ at each bond.
Periodic summation identifies the two occupation sums.\footnote{Formal declarations:
    \leanid{MPSTensor.parentInteraction_pvbs_two_apply},
    \leanid{MPSTensor.parentHamiltonianES_pvbs_decomposition}, and
    \leanid{MPSTensor.parentHamiltonianES_pvbs_gap_exact}.}

\section{Lower bound and attainment}
For $\mu\ge0$, the critical Hamiltonian $H_1$ and the coefficient $a(\mu)$
are nonnegative. Every configuration except the vacuum has at least one
particle. Thus for a vector with zero vacuum amplitude,
\begin{align}
    \langle v,H_\mu v\rangle & \ge c(\mu)\|v\|^2,
                             & \|H_\mu v\|        & \ge c(\mu)\|v\|.
    \label{eq:bn12_gap_bound}
\end{align}
In particular this holds on the orthogonal complement of the actual kernel,
since the vacuum is a zero-energy vector.

The uniform one-particle vector $W_N=\sum_j|0\cdots1_j\cdots0\rangle$
is annihilated by every critical local term and has particle number one.
Hence $H_\mu W_N=c(\mu)W_N$. For positive $\mu\ne1$, its eigenvalue is
positive, and symmetry makes $W_N$ orthogonal to the kernel. This nonzero
vector attains the bound, so the exact finite-ring gap is $c(\mu)$.

At $\mu=1$, the same vector joins the ground space. The vanishing of
$c(1)$ does not say that a fixed finite ring has zero gap above its enlarged
kernel. The separate Fourier-magnon argument below supplies nonzero low-energy
excitations as $N$ grows.

\section{Critical thermodynamic limit}
On the entire one-particle sector at $\mu=1$, the double-occupation
projector vanishes and the local interaction is the singlet exchange term.
The actual parent therefore acts by the normalized cycle Laplacian on the
particle amplitudes. The first Fourier character gives a nonzero vector
$\psi_N$ with
\begin{align}
    H_1\psi_N & =\eta_N\psi_N,
              & \eta_N         & =1-\cos(2\pi/N)>0\quad(N\ge2).
    \label{eq:bn12_critical_fourier}
\end{align}
Since the Hamiltonian is symmetric and $\eta_N>0$, this vector is
orthogonal to its entire kernel. As $\eta_N\to0$, any proposed positive
norm-gap bound on all sufficiently large rings is contradicted by this
family. This is the finite-periodic-volume criterion of thermodynamic
gaplessness, without constructing an infinite-volume operator or claiming
that any fixed finite ring has zero gap above its full kernel.\footnote{Formal
    declarations: \leanid{MPSTensor.parentHamiltonianES_pvbs_critical_fourierMagnon}
    and \leanid{MPSTensor.not_exists_eventual_uniform_gap_pvbs_critical}.}

\section{Remaining scope}
The coefficient and norm-gap theorems prove the positive-hopping
one-species finite-ring assertion. The general
multi-species thermodynamic conjecture, nonzero phases, and open-chain
spectral-gap estimates remain outside this statement. The open verdict
tracks that broader source scope, not a missing hypothesis in the finite
one-species result.
\end{document}
